

This problem highlights the strain invariance %that is
exhibited by the external application of the \Spherical{} transformation to the wrapped
elements. 
Originally presented by \cite{crisfield1999objectivity} for a 2-node element,
the analysis involves evaluating the curvature
\(\bm{\kappa} = \Twist \FrameBasis + \kappa_\alpha \FrameBasis_\alpha\)
in an element of length \(L = 1\) under
the following imposed rotations at the first and the last node, respectively:
\[
\bm{\Lambda}_1= \operatorname{Exp} \left(\begin{array}{c}
1 \\
-0.5 \\
\hphantom{-}0.25
\end{array}\right),
%
\quad \bm{\Lambda}_n= \operatorname{Exp}\left(\begin{array}{c}
-0.4 \\
\hphantom{-}0.7 \\
\hphantom{-}0.1
\end{array}\right).
\]
The current study extends the investigation by  
%In the test by 
\cite{jelenić1999geometrically}
to \(n\)-node elements %, only two-node elements were used.
with the interior nodes of higher-order elements subjected to 
identity rotations so that \(\bm{\Lambda}_a = \bm{1}\) when \(a \notin \{1,n\}\).

New curvatures
\(\bm{\kappa}^+ = \Twist^+ \FrameBasis + \kappa_\alpha^+\FrameBasis_\alpha\) are then determined 
in a coordinate system that results from a uniform rotation:
\[
\quad \bm{R} = \operatorname{Exp}\left(\begin{array}{c}
\hphantom{-}0.2 \\
\hphantom{-}1.2 \\
-0.5
\end{array}\right)
\]
so that the imposed rotations are given by \(\bm{\Lambda}_a^+ = \bm{R}\bm{\Lambda}_a\) at each node.

Table \ref{tab:objective} reports the components of \(\bm{\kappa}\) and \(\bm{\kappa}^+\) at the first integration
point.
Because the loading consists only of nodal rotations
imposed in a single increment 
and only strains are evaluated, the results for all three wrapped elements
are identical and reported only once in the table.
The results for two-node elements under the \texttt{None} and \Spherical{} 
transformations match the values %derived
in \cite{crisfield1999objectivity} exactly.

\begin{center}
\begin{table}[!h]%
\caption{%
Curvature components of Example~\ref{sec:frame-1000} for elements with \(n\) nodes.
All cases using the \Spherical{} isometry
show perfect agreement between \(\bm{\kappa}\) and \(\bm{\kappa}^+\).}
\label{tab:objective}

%
\begin{tabular*}{\textwidth}{@{\extracolsep\fill}cc|cc|cc|cc@{}}
\toprule
Isometry & \(n\) & \(\tau\) & \(\tau^+\) & \(\kappa_1\) & \(\kappa_1^+\) & \(\kappa_2\) & \(\kappa_2^+\) \\
\midrule %      k1          k*1         k2         k*2         k3           k*3
\texttt{None}    & 2 & -1.274645 & -1.263986  & 1.267559 & 1.313713 & -0.403500  & -0.337510 \\
\texttt{None}    & 3 & -2.121770 & -2.123805  & 0.951721 & 0.971454 & -0.463955  & -0.450500 \\
\texttt{None}    & 4 & -3.676150 & -3.657906  & 1.948050 & 2.025909 & -0.975779  & -0.919891 \\
\hline                                                                  
\texttt{Axis}    & 2 & -1.274236 & -1.242486  & 1.265924 & 1.316530 & -0.408787  & -0.373875 \\
\texttt{Axis}    & 3 & -2.122387 & -2.121660  & 0.951274 & 0.971569 & -0.463604  & -0.456998 \\
\texttt{Axis}    & 4 & -3.677140 & -3.647974  & 1.943930 & 2.015658 & -0.979285  & -0.959078 \\
\hline                                                                  
\texttt{SFIN}   & 2 & -1.263825 & -1.263825  & 1.271015 & 1.271015 & -0.422941  & -0.422941 \\
\texttt{SFIN}   & 3 & -2.121770 & -2.121770  & 0.951721 & 0.951721 & -0.463955  & -0.463955 \\
\texttt{SFIN}   & 4 & -3.676150 & -3.676150  & 1.948050 & 1.948050 & -0.975779  & -0.975779 \\
\hline
\end{tabular*}
\end{table}
\end{center}

\subparagraph{Discussion}

The agreement between the components of \(\bm{\kappa}\) and \(\bm{\kappa}^+\)
% \FCF{unclear what the following means} 
for cases using the \Spherical{} isometry
confirms that strain objectivity is restored for the wrapped elements.